% RESUMEN (IMRyD, ~300 palabras) + ABSTRACT en inglés.
\chapter*{Resumen}\addcontentsline{toc}{chapter}{Resumen}

La generación automática de informes de radiografía de tórax mediante modelos multimodales de
lenguaje (MLLM) es prometedora, pero su fiabilidad y la fidelidad de sus explicaciones visuales están
poco estudiadas, lo que limita su adopción clínica. Este trabajo evalúa si la explicación visual de
estos modelos (o su propia incertidumbre) permite estimar la fiabilidad de cada informe. Se
implementa un marco de explicabilidad y se valida el anclaje de la atención de forma objetiva contra
las cajas delimitadoras reales de NIH ChestX-ray8. El estudio de fiabilidad se articula como un diseño
\(2\times2\) que cruza dos generadores del estado del arte (LLaVA-Rad y MAIRA-2) con dos dominios
(Open-I, Estados Unidos; y PadChest-GR, España), etiquetando cada informe como correcto, alucinación u
omisión mediante tres fuentes de verdad independientes (diagnóstico estructurado MeSH, hallazgos
presentes/ausentes de radiólogo y juez basado en modelo de lenguaje). Los resultados muestran que la
atención está clínicamente anclada (lift global 3,03), pero la morfología de la explicación no predice
la fiabilidad. Este resultado se confirma incluso con las cajas de anclaje entrenadas del propio
modelo, que localiza sus alucinaciones de forma plausible. En cambio, la incertidumbre de la decisión, medida
como la probabilidad mínima de contenido, detecta las alucinaciones con AUC 0,90 (Open-I) y se
transfiere de forma robusta a otro modelo y a otro dominio (AUC 0,80--0,81), salvo en un generador
débil fuera de distribución. La señal es asimétrica: detecta las fabricaciones, pero no las omisiones,
que son errores confiados de sensibilidad. Una vez calibrada, la señal se traduce en una confianza por
informe con bandas de riesgo (``semáforo'') y en una demostración interactiva que integra informe,
anclaje y confianza. El hallazgo más relevante del trabajo es que, en los modelos y dominios
evaluados, la detección de alucinaciones mediante incertidumbre generativa resulta una señal
de fiabilidad más útil que la morfología de la explicación visual: es cuantificable, se transfiere
entre modelos y dominios, y es apta para el cribado clínico.

\vspace{0.4cm}\noindent\textbf{Palabras clave:} informes radiológicos, modelos multimodales,
explicabilidad, detección de alucinaciones, incertidumbre.

\chapter*{Abstract}\addcontentsline{toc}{chapter}{Abstract}

Automatic chest X-ray report generation with multimodal large language models (MLLM) is promising, but
the reliability of these models and the faithfulness of their visual explanations remain understudied,
hindering clinical adoption. This work assesses whether a model's visual explanation (or its own
uncertainty) can estimate the reliability of each report. An explainability framework is implemented
and the grounding of the model's attention is objectively validated against the real bounding boxes of
NIH ChestX-ray8. The reliability study follows a \(2\times2\) design crossing two state-of-the-art
generators (LLaVA-Rad and MAIRA-2) with two domains (Open-I, United States; and PadChest-GR, Spain),
labelling each report as correct, hallucination or omission through three independent ground-truth
sources (structured MeSH diagnosis, radiologist present/absent findings, and a language-model judge).
Results show that attention is clinically grounded (global lift 3.03), yet the morphology of the
explanation does not predict reliability. This finding is confirmed even with the model's own trained
grounding boxes, which plausibly localise its hallucinations. Instead, decision uncertainty, measured
as the minimum content-token probability, detects hallucinations with AUC 0.90 (Open-I) and transfers
robustly across model and domain (AUC 0.80--0.81), except for a weak generator out of distribution.
The signal is asymmetric: it detects fabrications but not omissions, which are confident sensitivity
errors. Once calibrated, the signal yields a per-report confidence with risk bands (a ``traffic
light'') and an interactive demo integrating report, grounding and confidence. The most relevant
finding of this work is that, across the models and domains evaluated, hallucination detection
via generative uncertainty is a more useful reliability signal than the morphology of the visual
explanation: it is quantifiable, transfers across models and domains, and is suitable for clinical
triage.

\vspace{0.4cm}\noindent\textbf{Keywords:} radiology reporting, multimodal models, explainability,
hallucination detection, uncertainty.
